\documentclass[a4paper,12pt]{article}
\usepackage[utf8]{inputenc}
\usepackage[T1]{fontenc}
\usepackage[ngerman]{babel}
\usepackage{amsmath, amssymb}
\usepackage{geometry}
\usepackage{tikz}
\usetikzlibrary{shapes.geometric, arrows.meta, positioning}
\usepackage{parskip}
\usepackage{titlesec}
\usepackage{enumitem}

\geometry{left=3cm, right=3cm, top=2.5cm, bottom=2.5cm}
\titleformat{\section}{\large\bfseries}{}{0em}{}[\titlerule]
\titleformat{\subsection}{\normalsize\bfseries}{}{0em}{}

\tikzset{
  box/.style={rectangle, draw, rounded corners=3pt, text width=4.8cm,
              align=center, minimum height=1.0cm, font=\small, inner sep=4pt},
  arrow/.style={-{Latex[length=2mm]}, thick},
}

\begin{document}

\begin{center}
  {\LARGE\bfseries Studierhinweise zu Lektion 4}\\[0.5em]
  {\large Lineare Algebra (Modul 61112)}
\end{center}

\vspace{0.8em}

Diese Lektion behandelt die \textbf{Determinante} -- eine der zentralen Funktionen der Linearen Algebra.
Die Determinante $\det(A) \in R$ einer quadratischen Matrix $A \in M_{n,n}(R)$ entscheidet über die
Invertierbarkeit von $A$ und wird in Lektion~5 zur Definition des charakteristischen Polynoms
benötigt. Die Lektion ist wichtig, aber zugänglich -- viele Determinanten rechnen zu üben hilft
schnell, ein intuitives Gefühl zu entwickeln.

% -----------------------------------------------------------------------
\section*{Struktur der Lektion 4}

\begin{center}
\begin{tikzpicture}[node distance=0.55cm]
  \node[box] (sym)  {4.1 Symmetrische Gruppen $S_n$\\Permutationen, Signatur};
  \node[box, below=0.5cm of sym] (det) {4.2 Determinante\\Leibniz-Formel, Charakterisierungssatz,\\Multiplikationssatz};
  \node[box, below=0.5cm of det] (min) {4.3 Minoren und Adjunkte\\Entwicklungssatz, Adjunktensatz};
  \node[box, below=0.5cm of min] (end) {4.4 Det.\ von Endomorphismen\\Basisunabhängigkeit};
  \draw[arrow] (sym) -- (det);
  \draw[arrow] (det) -- (min);
  \draw[arrow] (min) -- (end);
\end{tikzpicture}
\end{center}

\newpage

% -----------------------------------------------------------------------
\section*{Zielelement 4.1 -- Symmetrische Gruppen}

\subsection*{Lerninhalte}
\begin{center}
\begin{tikzpicture}[node distance=0.55cm]
  \node[box] (def) {Symmetrische Gruppe $S_n$\\$|S_n| = n!$ (Satz 4.1.3)};
  \node[box, below=0.5cm of def] (per) {Permutationen: Darstellung,\\Komposition, Träger};
  \node[box, below=0.5cm of per] (tra) {Transpositionen (4.1.14)\\jede Permutation = Produkt von Transpositionen};
  \node[box, below=0.5cm of tra] (sig) {Signatur $\mathrm{sgn}(\sigma) \in \{+1,-1\}$\\Satz von der Signatur (4.1.24)};
  \draw[arrow] (def) -- (per);
  \draw[arrow] (per) -- (tra);
  \draw[arrow] (tra) -- (sig);
\end{tikzpicture}
\end{center}

\subsection*{Lernziele}
\begin{itemize}
  \item Permutationen in Tabellenform lesen und Kompositionen ausrechnen.
  \item Jede Permutation als Produkt von Transpositionen darstellen.
  \item Die Signatur einer Permutation über Fehlstände oder Transpositionen bestimmen.
  \item Den Satz~4.1.24 kennen: $\mathrm{sgn}(\sigma \tau) = \mathrm{sgn}(\sigma)\cdot\mathrm{sgn}(\tau)$.
\end{itemize}

\subsection*{Selbstkontrollelement 4.1}
Berechnen Sie die Signatur der Permutation $\sigma = \begin{pmatrix}1&2&3&4\\3&4&2&1\end{pmatrix}$.

\newpage

% -----------------------------------------------------------------------
\section*{Zielelement 4.2 -- Determinante}

\subsection*{Lerninhalte}
\begin{center}
\begin{tikzpicture}[node distance=0.55cm]
  \node[box] (lei) {Leibniz-Formel (4.2.1)\\$\det(A) = \sum_{\sigma \in S_n} \mathrm{sgn}(\sigma) \prod_i a_{i,\sigma(i)}$};
  \node[box, below=0.5cm of lei] (drm) {Dreiecksmatrix: $\det = \prod a_{ii}$\\(4.2.13)};
  \node[box, below=0.5cm of drm] (cha) {Charakterisierungssatz (4.2.17)\\det eindeutig durch 3 Eigenschaften};
  \node[box, below=0.5cm of cha] (mul) {Multiplikationssatz (4.2.21)\\$\det(AB) = \det(A)\det(B)$};
  \draw[arrow] (lei) -- (drm);
  \draw[arrow] (drm) -- (cha);
  \draw[arrow] (cha) -- (mul);
\end{tikzpicture}
\end{center}

\subsection*{Lernziele}
\begin{itemize}
  \item Die Leibniz-Formel kennen und für kleine Matrizen ($2{\times}2$, $3{\times}3$) anwenden.
  \item Den Charakterisierungssatz (4.2.17) kennen: det ist die eindeutige alternierende multilineare Funktion mit $\det(I) = 1$.
  \item Den Multiplikationssatz anwenden: $\det(AB) = \det(A)\det(B)$, $\det(A^{-1}) = \det(A)^{-1}$.
  \item Die Determinante durch elementare Zeilenumformungen berechnen.
  \item $A$ invertierbar $\iff$ $\det(A) \neq 0$.
\end{itemize}

\subsection*{Selbstkontrollelement 4.2}
Berechnen Sie $\det\begin{pmatrix}2&1&0\\-1&3&2\\0&1&4\end{pmatrix}$ mithilfe von Zeilenumformungen.

\newpage

% -----------------------------------------------------------------------
\section*{Zielelement 4.3 -- Minoren, Entwicklungssatz und Adjunkte}

\subsection*{Lerninhalte}
\begin{center}
\begin{tikzpicture}[node distance=0.55cm]
  \node[box] (min) {Minor $\det(A_{ij})$\\gestrichene Zeile und Spalte};
  \node[box, below=0.5cm of min] (ent) {Entwicklungssatz (Laplace-Entwicklung)\\Entwicklung nach Zeile $i$ (4.3.7)};
  \node[box, below=0.5cm of ent] (adj) {Adjunkte $\mathrm{adj}(A)$\\(4.3.11)};
  \node[box, below=0.5cm of adj] (ads) {Adjunktensatz (4.3.14)\\$A \cdot \mathrm{adj}(A) = \det(A) \cdot I_n$};
  \draw[arrow] (min) -- (ent);
  \draw[arrow] (ent) -- (adj);
  \draw[arrow] (adj) -- (ads);
\end{tikzpicture}
\end{center}

\subsection*{Lernziele}
\begin{itemize}
  \item Minoren einer Matrix berechnen.
  \item Die Laplace-Entwicklung nach einer Zeile oder Spalte anwenden.
  \item Die Adjunkte einer Matrix bestimmen.
  \item Den Adjunktensatz kennen und zur Berechnung von $A^{-1}$ einsetzen: $A^{-1} = \frac{1}{\det(A)}\mathrm{adj}(A)$.
\end{itemize}

\subsection*{Selbstkontrollelement 4.3}
Berechnen Sie $A^{-1}$ für $A = \begin{pmatrix}1&2\\3&4\end{pmatrix}$ mithilfe des Adjunktensatzes.

\newpage

% -----------------------------------------------------------------------
\section*{Zielelement 4.4 -- Determinante von Endomorphismen}

\subsection*{Lerninhalte}
\begin{center}
\begin{tikzpicture}[node distance=0.55cm]
  \node[box] (def) {Definition $\det(\varphi)$ für $\varphi \in \mathrm{End}(V)$\\über Matrixdarstellung};
  \node[box, below=0.5cm of def] (unab){Basisunabhängigkeit:\\ähnliche Matrizen haben gleiche Det.};
  \draw[arrow] (def) -- (unab);
\end{tikzpicture}
\end{center}

\subsection*{Lernziele}
\begin{itemize}
  \item Die Determinante eines Endomorphismus als basisunabhängige Größe definieren.
  \item Verstehen, warum ähnliche Matrizen dieselbe Determinante haben (4.2.26).
  \item $\det(\varphi)$ an einem konkreten Endomorphismus berechnen.
\end{itemize}

\subsection*{Selbstkontrollelement 4.4}
Sei $\varphi: \mathbb{R}^2 \to \mathbb{R}^2$ die Drehung um $90°$. Berechnen Sie $\det(\varphi)$ und interpretieren Sie das Ergebnis geometrisch.

\end{document}
